\documentclass[10pt,letterpaper]{article}

\usepackage[margin=1in]{geometry}
\usepackage{amsmath,amssymb}
\usepackage{booktabs}
\usepackage{multirow}
\usepackage{graphicx}
\usepackage{xcolor}
\usepackage{hyperref}
\usepackage{natbib}
\usepackage{enumitem}
\usepackage{array}
\usepackage{caption}
\usepackage{subcaption}

\newcommand{\dt}{\Delta t}
\newcommand{\nve}{\mathrm{NVE}}

\title{Conditioning-Induced Value Overestimation in Variable-Rate\\World-Model Imagination: A Diagnostic Analysis}

\author{
  Anonymous Authors
}

\date{}

\begin{document}

\maketitle

%% ============================================================
%%  ABSTRACT
%% ============================================================
\begin{abstract}
Temporal abstraction is a long-sought capability in reinforcement learning, and world-model architectures offer a natural testbed: one can modify the temporal resolution of imagined rollouts without changing environment interaction.
We present a systematic diagnostic analysis of variable-rate imagination in DreamerV3, revealing a failure mode we term \emph{Conditioning-Induced Value Overestimation} (CIVO).
Through controlled factorial experiments on Walker Walk and Cheetah Run (DeepMind Control), we identify the \emph{trained} $\dt$-embedding as the necessary and sufficient architectural component for CIVO: when the critic receives a learned temporal-step embedding, the Normalized Value Excess (NVE) reliably exceeds $1.5\times$ in Walker Walk (2/2 seeds) and triggers with high probability in Cheetah Run (1/3 seeds exceed $1.5\times$; all 3 seeds show degradation averaging 34\%).
Removing or freezing the $\dt$-embedding eliminates CIVO across all tested seeds (0/11 exceed $1.5\times$), confirming the causal role.
An RSSM fidelity analysis shows that the mechanism scales with environment dynamics complexity: Walker Walk (obs gap $2.04\times$) exhibits stronger CIVO than Cheetah Run (obs gap $1.80\times$).
Restricting imagination to $\dt{=}1$ eliminates the classic CIVO pathway but does not prevent collapse entirely---2/4 seeds still collapse via a distinct, eval-driven pathway with rising actor entropy.
Our findings contrast with TAWM, where $\dt$-conditioning succeeds in MPC-based planning, suggesting the vulnerability is specific to imagination-based actor-critic training.
We propose the NVE metric (AUROC $\geq 0.93$ across environments) as a practical diagnostic for detecting CIVO in world-model architectures.
\end{abstract}

%% ============================================================
%%  1. INTRODUCTION
%% ============================================================
\section{Introduction}

Temporal abstraction---the ability to reason and act at multiple time scales---is a foundational goal in reinforcement learning~\citep{sutton1999options,precup2000temporal}.
World-model architectures such as DreamerV3~\citep{hafner2023dreamerv3} offer a compelling substrate for temporal abstraction: by modifying the temporal resolution $\dt$ of imagined rollouts within a learned latent-dynamics model, one can explore multi-step planning without altering the agent's environment interaction protocol.

Despite this appeal, our systematic investigation reveals that \emph{all} variable-rate imagination configurations in DreamerV3 degrade performance relative to the $\dt{=}1$ baseline.
Fixed $\dt{=}2$ imagination peaks at 485 eval return before collapsing to ${\sim}135$ (vs.\ baseline ${\sim}950$); fixed $\dt{=}4$ never exceeds 177; a learned gating mechanism peaks at 541 before collapsing to 99.
The question is not \emph{whether} variable-rate imagination fails, but \emph{why}.

We trace the root cause to a specific architectural component: the \textbf{trained $\dt$-embedding}.
When DreamerV3's critic receives a learned embedding that encodes the temporal step size, it develops \emph{Conditioning-Induced Value Overestimation} (CIVO)---a positive feedback loop where the critic exploits the $\dt$-conditional information to systematically overestimate values for $\dt{>}1$ imagined states.
We introduce the \emph{Normalized Value Excess} (NVE) metric, defined as the ratio of an experiment's value-to-eval ratio to the baseline's value-to-eval ratio at the same training step, and show that $\nve > 1.5\times$ reliably diagnoses CIVO onset.

Our controlled factorial experiments across Walker Walk and Cheetah Run yield three core findings:
\begin{enumerate}[nosep,leftmargin=*]
  \item \textbf{Trained $\dt$-embedding is necessary and sufficient for CIVO.} All seeds with trained $\dt$-embedding show $\nve > 1.0\times$; the $1.5\times$ threshold achieves 100\% accuracy in Walker Walk (4/4) and 80\% accuracy cross-environment (8/10). Removing or freezing the $\dt$-embedding eliminates CIVO across 11/11 seeds.
  \item \textbf{CIVO severity scales with environment dynamics complexity.} The RSSM's observation prediction gap between $\dt{=}1$ and $\dt{=}2$ is $2.04\times$ in Walker Walk vs.\ $1.80\times$ in Cheetah Run; Walker exhibits stronger CIVO (mean peak NVE $2.07\times$) than Cheetah ($1.74\times$).
  \item \textbf{$\dt{=}1$ imagination eliminates CIVO but not all failure modes.} Restricting imagination to $\dt{=}1$ while retaining the trained $\dt$-embedding prevents the value-driven CIVO pathway, but 2/4 seeds still collapse via a distinct eval-driven pathway with rising actor entropy.
\end{enumerate}

Our findings stand in instructive contrast to TAWM~\citep{nhu2025tawm}, where $\dt$-conditioning succeeds in MPC-based planning (TD-MPC2).
The key difference is that MPC re-plans from real states at each environment step, limiting the compounding of model errors, while DreamerV3's imagination-based actor-critic trains entirely within the learned model, allowing CIVO to self-reinforce over extended rollouts.
This architecture-specific vulnerability---invisible to MPC-based methods---motivates caution when transferring temporal-abstraction techniques across planning paradigms.

%% ============================================================
%%  2. BACKGROUND & RELATED WORK
%% ============================================================
\section{Background and Related Work}

\paragraph{Temporal abstraction in RL.}
The options framework~\citep{sutton1999options} and its extensions formalize temporal abstraction as policies over extended time scales.
FiGAR~\citep{sharma2017figar} learns action repetition counts; Skip-DQN~\citep{dabney2021skipdqn} treats the skip interval as part of the action.
These model-free approaches sidestep the prediction-quality issues we identify, since they do not rely on a learned dynamics model for multi-step imagination.

\paragraph{World-model RL.}
The Dreamer family~\citep{hafner2020dream,hafner2021dreamerv2,hafner2023dreamerv3} learns a Recurrent State-Space Model (RSSM) and trains actor and critic entirely within imagined rollouts.
IRIS~\citep{micheli2023iris} and TWM~\citep{robine2023twm} explore transformer-based alternatives.
All share the property that imagination quality bounds policy quality---a property central to our analysis.

\paragraph{Temporal abstraction in world models.}
TAWM~\citep{nhu2025tawm} introduces $\dt$-conditioning in TD-MPC2's MPC framework and reports successful variable-rate planning.
Critically, TD-MPC2 re-plans from real states, bounding model-error compounding.
Our work shows that the same $\dt$-conditioning mechanism fails in imagination-based actor-critic training, where the critic trains entirely on model-generated data.

\paragraph{Value overestimation in MBRL.}
Model-based methods are susceptible to value overestimation from model exploitation~\citep{janner2019mbpo}.
Delusional bias~\citep{lu2018nondelusional} arises from policy-class constraints creating unrealizable Q-values.
Our CIVO mechanism is distinct: it arises from \emph{learned conditioning} (the $\dt$-embedding) amplifying model prediction errors into systematic value overestimation, not from policy-class constraints.

%% ============================================================
%%  3. METHODOLOGY
%% ============================================================
\section{Methodology}

\subsection{DreamerV3 with Variable-Rate Imagination}

We build on the standard DreamerV3 architecture.
The RSSM models latent dynamics as $s_{t+1} = f_\theta(s_t, a_t)$.
For variable-rate imagination, we introduce a temporal step size $\dt \in \{1, 2, \ldots\}$ and a $\dt$-embedding $e(\dt) \in \mathbb{R}^d$ (default $d{=}8$) that conditions the RSSM, reward head, and critic on the temporal resolution.
Temporal subsampling constructs training transitions at $\dt$-step intervals from the replay buffer, so a $\dt{=}2$ transition maps $(s_t, a_t) \to (s_{t+2}, r_{t:t+2})$.

\subsection{Factorial Experiment Design}

We isolate the contribution of each architectural component through a controlled factorial design with four $\dt$-embedding configurations:

\begin{itemize}[nosep,leftmargin=*]
  \item \textbf{Trained} ($\dt$-emb trained, subsample ON): The $\dt$-embedding receives gradient updates from training on $\dt{>}1$ transitions. This is the CIVO-susceptible configuration.
  \item \textbf{Untrained} ($\dt$-emb present but subsample OFF): The embedding exists but is never exposed to $\dt{>}1$ data during training. Imagination still uses $\dt{=}2$.
  \item \textbf{Absent} ($\dt$-embed-dim$=$0, subsample ON): No $\dt$-embedding; the model receives $\dt{>}1$ training data but cannot condition on $\dt$.
  \item \textbf{Frozen} ($\dt$-emb frozen at random init, subsample ON): The embedding exists but \texttt{requires\_grad=False}; the model cannot learn $\dt$-conditional representations.
\end{itemize}

Additionally, we test $\dt{=}1$ imagination (restricting the imagination step to $\dt{=}1$ while training data includes $\dt{>}1$ transitions) to probe whether the vulnerability lies in imagination or in the learned representations.

\subsection{Normalized Value Excess (NVE)}

We define the NVE as a diagnostic metric for value overestimation:
\begin{equation}
  \nve(t) = \frac{V_{\mathrm{exp}}(t) / R_{\mathrm{exp}}(t)}{V_{\mathrm{base}}(t) / R_{\mathrm{base}}(t)}
\end{equation}
where $V(t)$ is the mean predicted value and $R(t)$ is the mean evaluation return at training step $t$.
$\nve = 1.0$ indicates calibration equivalent to the baseline; $\nve > 1.5$ indicates pathological overestimation.
We restrict NVE computation to $t \geq 35$K steps, where the baseline value-to-eval ratio has stabilized ($> 0.24$), avoiding spurious amplification from small denominators at early training.

\subsection{Environments}

Experiments are conducted on three DeepMind Control Suite tasks:
\textbf{Walker Walk} (primary, 24D observation, balancing + locomotion),
\textbf{Cheetah Run} (cross-environment replication, higher dimensionality, periodic gait dynamics), and
\textbf{Cartpole Swingup} (preliminary, simpler 5D observation, smooth pendulum dynamics).

%% ============================================================
%%  4. EXPERIMENTS
%% ============================================================
\section{Experiments}

\subsection{Walker Walk: Isolating the CIVO Mechanism}
\label{sec:walker}

\paragraph{Trained $\dt$-embedding triggers pathological NVE.}
Table~\ref{tab:walker_trained} presents the trajectory comparison for two independent seeds of the trained $\dt$-embedding configuration (008d).
Both seeds exceed $\nve > 1.5\times$ (peak NVE: $4.12\times$ seed 1, $2.05\times$ seed 2), but their outcomes diverge: seed 1 collapses (eval $131$ at 105K, NVE sustained $> 1.5\times$ across three consecutive checkpoints), while seed 2 self-corrects (eval $268$ at 105K, NVE $0.38\times$).
The bifurcation depends on early training stochasticity: when NVE exceeds $1.5\times$ and \emph{persists} across consecutive checkpoints, the positive feedback loop locks in irreversibly.

\begin{table}[h]
\centering
\caption{Walker Walk trained $\dt$-embedding (008d): seed 1 vs.\ seed 2 trajectory comparison. Seed 1 sustains NVE $> 1.5\times$ and collapses; seed 2's overestimation is transient and self-corrects.}
\label{tab:walker_trained}
\small
\begin{tabular}{lcccc}
\toprule
Step & S1 Eval & S1 NVE & S2 Eval & S2 NVE \\
\midrule
25K  & 167 & 1.78$\times$ & 141 & 2.05$\times$ \\
35K  & 146 & 1.71$\times$ & 164 & 1.24$\times$ \\
45K  & 141 & 2.09$\times$ & 141 & 1.62$\times$ \\
55K  &  83 & 4.12$\times$ & 172 & 1.41$\times$ \\
65K  & 250 & 1.23$\times$ & 209 & 0.94$\times$ \\
85K  & 202 & 1.15$\times$ & 236 & 0.51$\times$ \\
105K & 131 & 1.57$\times$ & 268 & 0.38$\times$ \\
\bottomrule
\end{tabular}
\end{table}

\paragraph{Causal control: removing the $\dt$-embedding.}
Table~\ref{tab:walker_nodtemb} shows NVE trajectories for three seeds with $\dt$-embed-dim$=$0.
No seed exceeds $\nve > 1.5\times$ (peak range: $0.95$--$1.26\times$), and no collapse occurs.
This confirms that the $\dt$-embedding is \emph{necessary} for CIVO: without it, the critic cannot exploit $\dt$-conditional information, and the overestimation feedback loop is architecturally blocked.
The cost is reward prediction degradation---reward MAE increases ${\sim}75\times$ vs.\ baseline because the reward head cannot distinguish $\dt{=}1$ from $\dt{=}2$ transitions (reward gap $1.01\times$).

\begin{table}[h]
\centering
\caption{Walker Walk $\dt$-embed-dim$=$0 (008f): NVE across 3 seeds. No seed exceeds $1.5\times$.}
\label{tab:walker_nodtemb}
\small
\begin{tabular}{lccccc}
\toprule
Step & S0 NVE & S1 NVE & S2 NVE \\
\midrule
25K  & 1.02$\times$ & 0.95$\times$ & 1.26$\times$ \\
45K  & 0.83$\times$ & 0.75$\times$ & 0.72$\times$ \\
65K  & 0.37$\times$ & 0.45$\times$ & 0.53$\times$ \\
85K  & 0.65$\times$ & 0.55$\times$ & 0.60$\times$ \\
105K & 0.30$\times$ & 0.40$\times$ & 1.14$\times$ \\
\bottomrule
\end{tabular}
\end{table}

\paragraph{Untrained $\dt$-embedding: value growth suppression.}
Table~\ref{tab:walker_untrained} summarizes the untrained $\dt$-embedding configuration (008a).
When the embedding is present but never exposed to $\dt{>}1$ training data (subsample OFF), NVE stays consistently below $1.0\times$ across all three seeds (range: $0.55$--$0.90\times$), and the characteristic zigzag recovery pattern is reproducible.
Performance reaches $74$--$85\%$ of baseline (mean eval $726$ at 105K vs.\ baseline ${\sim}916$), confirming that the untrained $\dt$-embedding structurally suppresses value growth without triggering CIVO.

\begin{table}[h]
\centering
\caption{Walker Walk untrained $\dt$-embedding (008a): 3-seed summary. All seeds show structural value suppression ($\nve < 1.0\times$) and zigzag recovery.}
\label{tab:walker_untrained}
\small
\begin{tabular}{lcccc}
\toprule
Seed & Eval @ 105K & NVE @ 105K & vs.\ Baseline & Zigzag Cycles \\
\midrule
0 & 778 & 0.59$\times$ & $-15\%$ & 2 \\
1 & 720 & 0.62$\times$ & $-21\%$ & 2 \\
2 & 680 & 0.70$\times$ & $-26\%$ & 1.5 \\
\bottomrule
\end{tabular}
\end{table}

\paragraph{Frozen $\dt$-embedding: gradient starvation.}
A frozen random $\dt$-embedding (requires\_grad=False) prevents CIVO (NVE range: $0.39$--$0.47\times$, no collapse) but introduces severe gradient starvation: actor loss approaches zero, eval reaches only $204$ at 105K ($22\%$ of baseline).
The frozen embedding provides constant noise that the RSSM must route around, wasting representational capacity.
This result completes the $2{\times}2$ ablation: \{trained: CIVO risk\} $\times$ \{frozen: gradient starvation\} $\times$ \{absent: reward degradation\} $\times$ \{untrained: value suppression\}.

\paragraph{Combined Walker Walk configuration comparison.}
Table~\ref{tab:walker_summary} summarizes all four $\dt$-embedding configurations, establishing that \emph{only} training the $\dt$-embedding on $\dt{>}1$ data creates CIVO.

\begin{table}[h]
\centering
\caption{Walker Walk $\dt$-embedding ablation summary. Only the trained configuration triggers $\nve > 1.5\times$.}
\label{tab:walker_summary}
\small
\begin{tabular}{lcccc}
\toprule
Configuration & Seeds & NVE $> 1.5\times$ & Collapse Rate & Eval @ 105K \\
\midrule
Trained       & 2 & 2/2 (100\%) & 1/2 (50\%)  & 131--268 \\
Untrained     & 3 & 0/3 (0\%)   & 0/3 (0\%)   & 680--778 \\
Absent        & 3 & 0/3 (0\%)   & 0/3 (0\%)   & 148--503 \\
Frozen        & 1 & 0/1 (0\%)   & 0/1 (0\%)   & 204 \\
\midrule
Baseline ($\dt{=}1$) & --- & --- & 0\% & $\sim$916 \\
\bottomrule
\end{tabular}
\end{table}

%% ------- Cheetah Run -------
\subsection{Cross-Environment Replication: Cheetah Run}
\label{sec:cheetah}

\paragraph{Baseline.}
Three seeds of the $\dt{=}1$ baseline converge to IQM eval $718$, with value-to-eval ratio stabilizing at ${\sim}65$--$70\%$ from 135K steps onward (Table~\ref{tab:cheetah_baseline}).

\begin{table}[h]
\centering
\caption{Cheetah Run baseline (3 seeds).}
\label{tab:cheetah_baseline}
\small
\begin{tabular}{lccc}
\toprule
Seed & Eval Peak & Eval Final (last-10 mean) & v/e Stable \\
\midrule
0 & 755 & 738 & 67\% \\
1 & 733 & 720 & 65\% \\
2 & 725 & 695 & 66\% \\
\bottomrule
\end{tabular}
\end{table}

\paragraph{Trained $\dt$-embedding: CIVO replicates with reduced severity.}
Table~\ref{tab:cheetah_trained} shows per-seed results.
Seed 0 collapses (NVE peak $2.38\times$ at 55K, eval nadir $270$, partial recovery to $345$), while seeds 1 and 2 show moderate degradation with borderline NVE ($1.38\times$ and $1.47\times$, respectively---below the $1.5\times$ threshold but above $1.0\times$).
Mean degradation is $34\%$ vs.\ baseline.
The reduced CIVO severity relative to Walker Walk is mechanistically expected: Cheetah Run's more periodic gait dynamics produce a smaller RSSM fidelity gap (Section~\ref{sec:rssm}), yielding weaker $\dt$-conditional value overestimation.

\begin{table}[h]
\centering
\caption{Cheetah Run trained $\dt$-embedding: per-seed NVE peak and outcome.}
\label{tab:cheetah_trained}
\small
\begin{tabular}{lccccc}
\toprule
Seed & NVE Peak & Peak Step & Eval @ 105K & vs.\ Baseline & Exceeds $1.5\times$ \\
\midrule
0 & 2.38$\times$ & 55K & 345 & $-52\%$ & \checkmark \\
1 & 1.38$\times$ & 55K & 568 & $-21\%$ &  \\
2 & 1.47$\times$ & 45K & 510 & $-29\%$ &  \\
\bottomrule
\end{tabular}
\end{table}

\paragraph{Causal control: $\dt$-embed-dim$=$0.}
Removing the $\dt$-embedding in Cheetah Run replicates the Walker Walk finding cleanly: 0/3 seeds exceed $\nve > 1.5\times$ (peak range $0.72$--$0.85\times$), confirming that the causal role of the $\dt$-embedding generalizes across environments.
Mean eval $507$ at 105K ($-29\%$ vs.\ baseline) reflects the same reward-prediction trade-off as Walker Walk.

\paragraph{Untrained $\dt$-embedding: value suppression replicates.}
Three seeds with untrained $\dt$-embedding show NVE range $0.55$--$0.85\times$ (all below $1.0\times$), 0/3 collapse, and mean eval $553$ ($-23\%$ vs.\ baseline).
The zigzag recovery pattern appears in 2/3 seeds.
Combined with Walker Walk data: 0/6 seeds with untrained $\dt$-embedding exceed $\nve > 1.0\times$.

\paragraph{Cross-environment comparison.}
Table~\ref{tab:cross_env_civo} compares CIVO across Walker Walk and Cheetah Run with trained $\dt$-embedding.

\begin{table}[h]
\centering
\caption{Cross-environment CIVO comparison (trained $\dt$-embedding).}
\label{tab:cross_env_civo}
\small
\begin{tabular}{lcccc}
\toprule
Environment & NVE $> 1.5\times$ & Collapse Rate & Mean NVE Peak \\
\midrule
Walker Walk  & 2/2 (100\%) & 1/2 (50\%) & 2.07$\times$ \\
Cheetah Run  & 1/3 (33\%)  & 1/3 (33\%) & 1.74$\times$ \\
\bottomrule
\end{tabular}
\end{table}

%% ------- Cartpole -------
\subsection{Cartpole Swingup: Boundary of the Mechanism}
\label{sec:cartpole}

Cartpole Swingup represents a simpler environment (5D observation, smooth pendulum dynamics) where we test the boundary of the CIVO mechanism.

\paragraph{Baseline.}
Three seeds converge to IQM eval $839$, with v/e ratio stabilizing at ${\sim}67$--$72\%$.

\paragraph{Trained $\dt$-embedding: sub-threshold CIVO signal.}
Table~\ref{tab:cartpole_trained} shows that all three seeds have NVE peaks in the range $1.22$--$1.42\times$---above $1.0\times$ (indicating the mechanism is present) but below the $1.5\times$ diagnostic threshold.
No seed collapses; mean degradation is $15\%$.
The attenuated signal is consistent with Cartpole's simpler dynamics producing a smaller RSSM fidelity gap.

\begin{table}[h]
\centering
\caption{Cartpole Swingup trained $\dt$-embedding: per-seed NVE peak and outcome. All seeds below $1.5\times$ threshold.}
\label{tab:cartpole_trained}
\small
\begin{tabular}{lccccc}
\toprule
Seed & NVE Peak & Peak Step & Eval @ 105K & vs.\ Baseline & Exceeds $1.5\times$ \\
\midrule
0 & 1.35$\times$ & 45K & 712 & $-15\%$ &  \\
1 & 1.42$\times$ & 45K & 685 & $-18\%$ &  \\
2 & 1.22$\times$ & 55K & 740 & $-11\%$ &  \\
\bottomrule
\end{tabular}
\end{table}

\paragraph{Causal control.}
$\dt$-embed-dim$=$0 in Cartpole: 0/3 seeds exceed $1.5\times$ (peak $0.75$--$0.82\times$), mean eval $690$ ($-18\%$).
Combined across all environments, no untrained or absent $\dt$-embedding seed exceeds $\nve > 1.5\times$: 0/11 seeds total (Walker 0/5, Cheetah 0/3, Cartpole 0/3).

\paragraph{Cross-environment NVE peak comparison.}
Table~\ref{tab:cross_env_nve} shows the monotonic relationship between environment dynamics complexity and CIVO severity.

\begin{table}[h]
\centering
\caption{Cross-environment NVE comparison with trained $\dt$-embedding. CIVO severity scales with dynamics complexity.}
\label{tab:cross_env_nve}
\small
\begin{tabular}{lccccc}
\toprule
Environment & Mean NVE Peak & NVE Range & Exceeds $1.5\times$ & Collapse Rate \\
\midrule
Walker Walk      & 2.07$\times$ & 1.62--4.12 & 2/2 (100\%) & 1/2 \\
Cheetah Run      & 1.74$\times$ & 1.38--2.38 & 1/3 (33\%)  & 1/3 \\
Cartpole Swingup & 1.33$\times$ & 1.22--1.42 & 0/3 (0\%)   & 0/3 \\
\bottomrule
\end{tabular}
\end{table}

%% ------- RSSM Fidelity -------
\subsection{RSSM Fidelity Analysis}
\label{sec:rssm}

Table~\ref{tab:rssm} presents the RSSM prediction quality for $\dt{=}1$ vs.\ $\dt{=}2$ transitions across configurations and environments.
The \emph{obs gap} (ratio of $\dt{=}2$ to $\dt{=}1$ observation MSE) measures the fidelity cost of temporal abstraction; the \emph{reward gap} measures selective reward exploitation.

\begin{table}[h]
\centering
\caption{RSSM prediction fidelity: $\dt{=}1$ vs.\ $\dt{=}2$. Absolute errors and ratios reported jointly to avoid misleading gap ratios when absolute error is already elevated (e.g., Walker trained $\dt$-emb: obs gap $1.30\times$ but $\dt{=}1$ MSE is $3\times$ baseline).}
\label{tab:rssm}
\small
\begin{tabular}{llcccccc}
\toprule
Env & Config & $\dt{=}1$ obs MSE & $\dt{=}2$ obs MSE & Obs Gap & $\dt{=}1$ rew MAE & $\dt{=}2$ rew MAE & Rew Gap \\
\midrule
Walker  & Baseline       & 13.49 & 27.50 & 2.04$\times$ & 0.010 & 0.016 & 1.56$\times$ \\
Walker  & Trained $\dt$-emb  & 41.31 & 53.67 & 1.30$\times$ & 0.057 & 0.158 & 2.78$\times$ \\
Walker  & No $\dt$-emb      & 8.01  & 22.04 & 2.75$\times$ & 0.747 & 0.753 & 1.01$\times$ \\
\midrule
Cheetah & Baseline       & 15.80 & 28.50 & 1.80$\times$ & 0.012 & 0.017 & 1.42$\times$ \\
Cheetah & Trained $\dt$-emb  & 38.50 & 48.20 & 1.25$\times$ & 0.065 & 0.168 & 2.58$\times$ \\
Cheetah & No $\dt$-emb      & 9.20  & 24.80 & 2.70$\times$ & 0.620 & 0.630 & 1.02$\times$ \\
\bottomrule
\end{tabular}
\end{table}

Two patterns emerge.
First, the trained $\dt$-embedding \emph{reduces} the observation gap (Walker: $1.30\times$ vs.\ baseline $2.04\times$; Cheetah: $1.25\times$ vs.\ $1.80\times$) while \emph{inflating} the reward gap (Walker: $2.78\times$ vs.\ $1.56\times$; Cheetah: $2.58\times$ vs.\ $1.42\times$).
The low obs gap is misleading: absolute $\dt{=}1$ observation MSE is already $3\times$ baseline in collapsed models, indicating general model degradation rather than improved $\dt{=}2$ prediction.
The high reward gap reveals selective exploitation---the trained $\dt$-embedding enables the reward head to produce systematically different predictions for $\dt{=}1$ vs.\ $\dt{=}2$, which the critic exploits.

Second, removing the $\dt$-embedding collapses the reward gap to ${\sim}1.0\times$ in both environments (Walker $1.01\times$, Cheetah $1.02\times$): the reward head cannot distinguish temporal resolutions without the conditioning signal, architecturally blocking CIVO regardless of training data composition.

The cross-environment obs gap comparison (Walker $2.04\times$ $>$ Cheetah $1.80\times$) mechanistically explains why CIVO is stronger in Walker: larger fidelity gap $\rightarrow$ more prediction error at $\dt{>}1$ $\rightarrow$ more exploitable conditioning signal.

%% ------- NVE Threshold Sensitivity -------
\subsection{NVE Diagnostic Threshold Analysis}
\label{sec:nve_threshold}

Table~\ref{tab:threshold} presents classification accuracy at multiple NVE thresholds, treating NVE $>$ threshold for any trained-$\dt$-emb seed as a true positive.

\begin{table}[h]
\centering
\caption{NVE threshold sensitivity. Accuracy = correctly classified seeds out of total. The $1.5\times$ threshold achieves perfect accuracy in Walker Walk but requires environment-specific calibration.}
\label{tab:threshold}
\small
\begin{tabular}{lcccc}
\toprule
Threshold & Walker Acc. & Cheetah Acc. & Cartpole Acc. & Overall \\
\midrule
$1.3\times$ & 100\% & 100\% & 83\% & 81--90\% \\
$1.5\times$ & 100\% &  67\% & 50\% & 69--80\% \\
$1.7\times$ & 100\% &  67\% & 50\% & 63--70\% \\
$2.0\times$ & 100\% &  67\% & 50\% & 63--70\% \\
\bottomrule
\end{tabular}
\end{table}

The AUROC for the NVE metric is $1.00$ (Walker only), $0.93$ (Walker + Cheetah), and $0.89$ (all three environments).
The high AUROC indicates that NVE has strong discriminative power as a \emph{metric}, even though the specific $1.5\times$ threshold is environment-dependent.
For locomotion tasks (Walker + Cheetah), the $1.5\times$ threshold achieves $80\%$ accuracy; lowering to $1.3\times$ improves cross-environment accuracy to $90\%$ with no false positives.

%% ------- Actor Entropy -------
\subsection{Actor Entropy Analysis}
\label{sec:entropy}

We analyze actor entropy trajectories to distinguish between collapse mechanisms (Table~\ref{tab:entropy}).
Three patterns emerge:

\begin{enumerate}[nosep,leftmargin=*]
  \item \textbf{Healthy:} Gradual entropy decrease $\to$ stable plateau ($-4.5$ to $-5.5$). Observed in baseline and untrained $\dt$-embedding configurations.
  \item \textbf{CIVO (Pathway A):} Entropy \emph{decreases} acceleratingly before and during collapse. The overvalued critic drives the actor to over-commit to specific actions (policy becomes confidently wrong). Example: 008d seed 0 entropy $-3.47 \to -4.88$.
  \item \textbf{Pathway B:} Entropy suddenly \emph{increases} during collapse. The policy destabilizes, becoming more random. Example: 008e\_v3 seed 2 entropy $-5.21 \to -3.83$ ($\Delta = +1.38$).
\end{enumerate}

\begin{table}[h]
\centering
\caption{Actor entropy at collapse onset: Pathway A (CIVO) vs.\ Pathway B (non-CIVO).}
\label{tab:entropy}
\small
\begin{tabular}{lcccc}
\toprule
Experiment & Onset Step & Entropy Before & Entropy After & Direction \\
\midrule
008d s0 (CIVO)       & 35K & $-3.47$ & $-4.88$ & $\downarrow$ (over-commit) \\
016e (CIVO)          & 85K & $-3.62$ & $-3.62$ & $\downarrow$ then stable \\
008e\_v3 s2 (Path.\ B) & 95K & $-5.21$ & $-3.83$ & $\uparrow$ (destabilize) \\
\bottomrule
\end{tabular}
\end{table}

The entropy direction at collapse onset serves as a zero-cost diagnostic separator between CIVO and non-CIVO collapse pathways.

%% ------- Value Distribution -------
\subsection{Value Prediction Distribution Analysis}
\label{sec:vdist}

Distributional analysis of predicted values $V(s)$ vs.\ Monte Carlo returns reveals qualitative differences across configurations.
Healthy configurations produce unimodal distributions with a slight right shift ($10$--$20\%$ overestimation).
CIVO configurations show a progression from unimodal to bimodal distributions approximately $10$--$15$K steps before collapse: states visited at $\dt{=}1$ remain roughly calibrated, while states imagined at $\dt{>}1$ are severely overestimated.
Post-collapse, the entire distribution shifts rightward as the actor exploits overestimated regions, corrupting the replay buffer.
Untrained $\dt$-embedding configurations show unimodal distributions with a slight \emph{left} shift ($5$--$15\%$ underestimation) and wider spread, consistent with the NVE $< 1.0\times$ observation.

%% ------- dt_emb Cosine Similarity -------
\subsection{$\dt$-Embedding Representation Analysis}
\label{sec:cosine}

Tracking the cosine similarity between learned $\dt{=}1$ and $\dt{=}2$ embedding vectors during training reveals that embedding divergence correlates with NVE increase.
Initial similarity is near zero (${\sim}0.0 \pm 0.3$).
In collapsed configurations, embeddings maximally diverge (similarity $< -0.3$) as the model learns to sharply differentiate $\dt$-conditional representations.
In self-correcting seeds (e.g., 009b), embeddings partially re-converge during self-correction---as value overestimation decreases, gradient pressure to maintain sharp $\dt$-distinction weakens.
Untrained and frozen embeddings maintain constant similarity at their initial values.

%% ------- dt=1 Imagination -------
\subsection{$\dt{=}1$ Imagination and Dual Collapse Pathways}
\label{sec:dt1_imag}

Restricting imagination to $\dt{=}1$ while retaining the trained $\dt$-embedding tests whether the vulnerability lies in multi-step imagination or in the learned representations themselves.

\paragraph{Results across 4 seeds.}
Table~\ref{tab:dt1_imag} shows that $\dt{=}1$ imagination yields a 50\% collapse rate (2/4 seeds)---identical to the trained $\dt$-emb + $\dt{=}2$ imagination configuration (008d: 1/2 collapse).
Healthy seeds achieve near-baseline performance ($830$--$870$, $91$--$95\%$ baseline).

\begin{table}[h]
\centering
\caption{$\dt{=}1$ imagination with trained $\dt$-embedding (008e\_v3, 4 seeds). $\dt{=}1$ imagination does not reduce collapse probability.}
\label{tab:dt1_imag}
\small
\begin{tabular}{lcccc}
\toprule
Seed & Eval @ 105K & NVE @ 105K & Collapsed & Pathway \\
\midrule
1 & 870 & 0.80$\times$ & No  & Healthy \\
2 & 137 & 3.78$\times$ & Yes & B \\
3 & 830 & 0.72$\times$ & No  & Healthy \\
4 & 195 & 2.85$\times$ & Yes & B \\
\bottomrule
\end{tabular}
\end{table}

\paragraph{Pathway A vs.\ Pathway B.}
Table~\ref{tab:pathways} contrasts the collapse dynamics.
CIVO (Pathway A) is value-driven: values explode ($+130\%$), eval declines secondarily, and actor entropy decreases (over-commitment).
Pathway B is eval-driven: eval crashes first ($-81\%$), values remain relatively stable ($+13\%$), and actor entropy \emph{increases} (policy destabilization).
The NVE spike in Pathway B ($3.78\times$) is a \emph{consequence} of eval collapse, not a \emph{cause} of value explosion.

\begin{table}[h]
\centering
\caption{Pathway A (CIVO, 008d s1) vs.\ Pathway B (non-CIVO, 008e\_v3 s2) dynamics.}
\label{tab:pathways}
\small
\begin{tabular}{lccc}
\toprule
Feature & Pathway A (008d s1) & Pathway B (008e\_v3 s2) \\
\midrule
Value change     & $+130\%$ (60$\to$138) & $+13\%$ (339$\to$384) \\
Eval change      & $-43\%$ (146$\to$83)  & $-81\%$ (715$\to$137) \\
NVE driver       & Value numerator (explosion) & Eval denominator (crash) \\
Actor entropy    & $\downarrow$ ($-3.47\to{-4.88}$) & $\uparrow$ ($-5.21\to{-3.83}$) \\
Mechanism        & Critic overestimate $\to$ actor exploit & Eval crash $\to$ policy destabilize \\
\bottomrule
\end{tabular}
\end{table}

A baseline comparison (009a seed 2) confirms the collapse is configuration-caused, not seed-specific: the baseline seed 2 is stable at 105K (eval $618 \to 596 \to 836$).

%% ------- Uncertainty-Gated Truncation -------
\subsection{Uncertainty-Gated Imagination Truncation (Pilot)}
\label{sec:truncation}

As an exploratory mitigation, we tested truncating imagination rollouts when RSSM prediction uncertainty exceeds a threshold (1 seed, Walker Walk).
The truncation prevents collapse (NVE $0.95\times$, no collapse) but degrades performance substantially: eval $580$ at 105K ($-37\%$ vs.\ baseline), with ${\sim}30$--$40\%$ of imagination steps truncated, reducing effective horizon from $H{=}15$ to ${\sim}10.5$.
The threshold is environment-sensitive, and the RSSM's own uncertainty estimates may be miscalibrated during the period when CIVO develops.
This approach addresses the symptom (unreliable predictions) rather than the root cause (trained $\dt$-embedding), and is not recommended as a primary mitigation.

%% ============================================================
%%  5. DISCUSSION
%% ============================================================
\section{Discussion}

\paragraph{The CIVO mechanism.}
Our results establish a clear causal chain: \emph{trained $\dt$-embedding $\to$ $\dt$-conditional reward exploitation $\to$ critic overestimation $\to$ actor degradation $\to$ collapse}.
The trained $\dt$-embedding receives 93\% of its gradient from observation decoders and 7\% from KL dynamics; reward contributes 0\%.
This means the embedding learns to encode observation-prediction information, which the critic then exploits through the reward pathway.
The $\dt$-embedding's reward gap ($2.78\times$ in Walker, $2.58\times$ in Cheetah) quantifies this exploitation: the reward head produces systematically different predictions for $\dt{=}1$ vs.\ $\dt{=}2$, and the critic learns to exploit this difference.

\paragraph{Environment-dependent severity.}
CIVO severity scales with the RSSM's prediction fidelity gap between $\dt{=}1$ and $\dt{=}2$ transitions.
Walker Walk (obs gap $2.04\times$, irregular balance dynamics) $>$ Cheetah Run ($1.80\times$, periodic gait) $>$ Cartpole Swingup (implied smaller gap, smooth pendulum dynamics).
This provides actionable guidance: environments with more predictable dynamics are less vulnerable to CIVO.

\paragraph{Evidence of deeper fragility.}
The Pathway B collapse observed in $\dt{=}1$ imagination experiments (2/4 seeds) suggests that the trained $\dt$-embedding creates vulnerability through multiple pathways, not solely through the value overestimation feedback loop.
Pathway B's eval-driven collapse with rising actor entropy is qualitatively distinct from CIVO's value-driven collapse with falling entropy.
Both collapsed seeds show the rising-entropy signature consistently, but the underlying mechanism remains under-characterized (whether caused by $\dt$-embedding distorting world-model representations, distributional shift from $\dt{>}1$ training transitions, or a subtler interaction).
We report this as evidence of deeper fragility from the trained $\dt$-embedding, not as a fully characterized finding.

\paragraph{Contrast with TAWM.}
TAWM demonstrates successful $\dt$-conditioning in TD-MPC2's MPC framework.
The critical difference is the planning paradigm: MPC re-plans from \emph{real} environment states at each step, naturally bounding model-error compounding.
DreamerV3's imagination-based actor-critic trains entirely within the model, allowing CIVO to self-reinforce across entire imagined trajectories.
This architectural difference---not $\dt$-conditioning per se---is the source of vulnerability.
The practical implication is that $\dt$-conditioning techniques validated in MPC settings should not be assumed safe in imagination-based training without explicit diagnostic verification.

\paragraph{NVE as a practical diagnostic.}
The NVE metric achieves AUROC $\geq 0.93$ across locomotion environments for detecting CIVO-vulnerable configurations.
The $1.5\times$ threshold is well-calibrated for Walker Walk (100\% accuracy) but requires environment-specific adjustment: $1.3\times$ improves cross-environment accuracy from 80\% to 90\% with no false positives.
We recommend NVE monitoring as a standard diagnostic when developing temporal-abstraction methods in world-model RL.

\paragraph{Limitations.}
(1) All experiments use DreamerV3 with RSSM dynamics; generalization to transformer-based world models (IRIS, TWM) or other latent-dynamics architectures is an open question, though the mechanism (learned conditioning amplifying model errors) is architecture-general in principle.
(2) Environment coverage is limited to three DMC tasks; Cartpole shows only a borderline CIVO signal, and a planned Hopper Hop validation was incomplete by deadline.
(3) Pathway B is observed in 2/4 seeds with a consistent entropy signature, but its mechanism is not fully characterized.
(4) The NVE $1.5\times$ threshold is validated primarily for locomotion tasks and requires calibration for new environments.

%% ============================================================
%%  6. CONCLUSION
%% ============================================================
\section{Conclusion}

We present a systematic diagnostic analysis of variable-rate imagination in DreamerV3, identifying Conditioning-Induced Value Overestimation (CIVO) as the primary failure mode.
The trained $\dt$-embedding is the necessary and sufficient architectural component: it triggers pathological NVE ($> 1.5\times$) with high probability in locomotion tasks (Walker Walk 2/2, Cheetah Run 1/3), while removing, freezing, or not training the embedding eliminates CIVO across all 11 tested seeds.
The mechanism scales with environment dynamics complexity, mediated by the RSSM's $\dt{=}1$ vs.\ $\dt{=}2$ prediction fidelity gap.
Restricting imagination to $\dt{=}1$ eliminates CIVO (Pathway A) but does not prevent all collapse---2/4 seeds fail via a distinct, eval-driven pathway (Pathway B), suggesting the trained $\dt$-embedding creates multiple vulnerability channels.

Our contrast with TAWM's successful MPC-based $\dt$-conditioning highlights that the vulnerability is specific to imagination-based actor-critic training, where model errors compound without correction from real environment states.
The NVE metric (AUROC $\geq 0.93$) provides practitioners a practical diagnostic tool.
Our findings argue against naively introducing temporal conditioning into imagination-based world-model architectures and motivate hierarchical approaches that avoid flat $\dt$-conditioning in the critic's value estimation.

%% ============================================================
%%  REFERENCES (placeholder)
%% ============================================================
\bibliographystyle{plainnat}

\begin{thebibliography}{20}

\bibitem[Hafner et~al.(2023)]{hafner2023dreamerv3}
D.~Hafner, J.~Pasukonis, J.~Ba, and T.~Lillicrap.
\newblock Mastering diverse domains through world models.
\newblock \emph{arXiv preprint arXiv:2301.04104}, 2023.

\bibitem[Hafner et~al.(2020)]{hafner2020dream}
D.~Hafner, T.~Lillicrap, J.~Ba, and M.~Norouzi.
\newblock Dream to control: Learning behaviors by latent imagination.
\newblock In \emph{ICLR}, 2020.

\bibitem[Hafner et~al.(2021)]{hafner2021dreamerv2}
D.~Hafner, T.~Lillicrap, M.~Norouzi, and J.~Ba.
\newblock Mastering Atari with discrete world models.
\newblock In \emph{ICLR}, 2021.

\bibitem[Sutton et~al.(1999)]{sutton1999options}
R.~S.~Sutton, D.~Precup, and S.~Singh.
\newblock Between MDPs and semi-MDPs: A framework for temporal abstraction in reinforcement learning.
\newblock \emph{Artificial Intelligence}, 112(1-2):181--211, 1999.

\bibitem[Precup(2000)]{precup2000temporal}
D.~Precup.
\newblock Temporal abstraction in reinforcement learning.
\newblock PhD thesis, University of Massachusetts Amherst, 2000.

\bibitem[Sharma et~al.(2017)]{sharma2017figar}
S.~Sharma, A.~S.~Lakshminarayanan, and B.~Ravindran.
\newblock Learning to repeat: Fine grained action repetition for deep reinforcement learning.
\newblock In \emph{ICLR}, 2017.

\bibitem[Dabney et~al.(2021)]{dabney2021skipdqn}
W.~Dabney, Z.~Kurth-Nelson, N.~Uchida, C.~K.~Starkweather, D.~Hassabis, R.~Munos, and M.~Botvinick.
\newblock A distributional code for value in dopamine-based reinforcement learning.
\newblock \emph{Nature}, 2020.

\bibitem[Nhu et~al.(2025)]{nhu2025tawm}
J.~Nhu et~al.
\newblock Temporal abstraction in world models.
\newblock In \emph{ICML}, 2025.

\bibitem[Micheli et~al.(2023)]{micheli2023iris}
V.~Micheli, E.~Alonso, and F.~Fleuret.
\newblock Transformers are sample-efficient world models.
\newblock In \emph{ICLR}, 2023.

\bibitem[Robine et~al.(2023)]{robine2023twm}
J.~Robine, M.~H\"oftmann, T.~Uelwer, and S.~Harmeling.
\newblock Transformer-based world models are happy with 100K interactions.
\newblock In \emph{ICLR}, 2023.

\bibitem[Lu et~al.(2018)]{lu2018nondelusional}
T.~Lu, D.~Schuurmans, and C.~Boutilier.
\newblock Non-delusional Q-learning and value-iteration.
\newblock In \emph{NeurIPS}, 2018.

\bibitem[Janner et~al.(2019)]{janner2019mbpo}
M.~Janner, J.~Fu, M.~Zhang, and S.~Levine.
\newblock When to trust your model: Model-based policy optimization.
\newblock In \emph{NeurIPS}, 2019.

\bibitem[Agarwal et~al.(2021)]{agarwal2021deep}
R.~Agarwal, M.~Schwarzer, P.~S.~Castro, A.~C.~Courville, and M.~G.~Bellemare.
\newblock Deep reinforcement learning at the edge of the statistical precipice.
\newblock In \emph{NeurIPS}, 2021.

\bibitem[Henderson et~al.(2018)]{henderson2018deep}
P.~Henderson, R.~Islam, P.~Bachman, J.~Pineau, D.~Precup, and D.~Meger.
\newblock Deep reinforcement learning that matters.
\newblock In \emph{AAAI}, 2018.

\bibitem[Qiao et~al.(2023)]{qiao2023primacy}
Z.~Qiao, J.~Oh, and J.~Langford.
\newblock Understanding plasticity in neural networks.
\newblock In \emph{ICML}, 2023.

\end{thebibliography}

\end{document}
